\section{From heights to finiteness}\label{fin:section}

We first descend the height estimate to an arithmetic estimate on the
original coarse curve.  We then identify the entire unpolarized
elliptic-square locus with the locus to which the counting theorem
applies. The two companion branch theorems enter only in the final
proof of the full Zilber--Pink theorem.

\subsection{The discriminant and the Galois orbit}

For an algebraic point $s$ with
$\End_{\Qbar}(A_s)\otimes\Q\simeq\Mat_2(\Q)$, put
\[
 \cO_s=\End_{\Qbar}(A_s),\qquad D_s=\cO_s\otimes\Q.
\]
Throughout this section, the discriminant of this order means
\begin{equation}\label{fin:disc-definition}
 \Delta_s=
 \left|\det\left(\Tr_{D_s/\Q}(e_i e_j)\right)_{1\leq i,j\leq4}\right|,
\end{equation}
where $(e_i)$ is any integral basis of $\cO_s$ and
$\Tr_{D_s/\Q}(x)$ is the trace of left multiplication by $x$ on $D_s$.
In particular, the discriminant is that of the full order, and includes
its conductor.  It is a positive integer independent of the basis.

\begin{proposition}[Orbit estimate]\label{fin:orbit}
Let $C\subset\A_2$ be a curve as in \Cref{intro:squares}.
There are a number field $K$ defining $C$ and constants $c,\delta>0$
such that every $s\in\Sigma_{E^2}(C)$ satisfies
\begin{equation}\label{fin:orbit-bound}
 \#\bigl(\operatorname{Aut}(\C/K)\cdot s\bigr)
   =[K(s):K]\ \geq\ c\Delta_s^\delta.
\end{equation}
The constants are uniform over the whole indicated locus.
\end{proposition}

\begin{proof}
If this locus is finite, the assertion follows by taking the minimum
of finitely many positive ratios; the empty case is vacuous.
Otherwise, make the fixed choices in \Cref{geo:setup,inc:markings}.
All the resulting covers, auxiliary primes, and the center are fixed.
Enlarge the fixed number field $K$ to contain their fields of definition.
After removing a finite set $F\subset C(\Qbar)$, their composite gives
a finite cover $q:S\longrightarrow C\setminus F$ of fixed degree $N$.
Further fixed deletions in the preceding sections are included in $F$
by removing their images.  We can restrict to the locus over which the
cover has constant degree, and include its remaining exceptional points
in $F$ as well.

For $s\in\Sigma_{E^2}(C)\setminus F$, choose the lift $t\in S(\Qbar)$
with the two markings used in \Cref{nid:height}.  Write
$r_s=[K(s):K]$ and $d_t=\max\{2,[K(t):K]\}$.
A closed point in a fiber of a finite map of degree $N$ has residue
degree at most $N$.  Thus, for every such chosen lift,
\begin{equation}\label{fin:degree-descent}
 [K(t):K]\leq N r_s,\qquad d_t\leq 2N r_s.
\end{equation}
This bound uses the degree of the fixed cover even though the chosen
point in its fiber depends on $s$.

The height bound in \Cref{nid:height} and the full endomorphism estimate
in \Cref{ht:endomorphisms} give fixed positive constants $C_1,C_2,u,v$
such that
\begin{equation}\label{fin:height-to-disc}
 h(t)\leq C_1 d_t^u,\qquad
 \Delta_s=\Delta_t\leq C_2\bigl(d_t h(t)\bigr)^v.
\end{equation}
Here forgetting a level structure leaves the geometric principally
polarized abelian variety unchanged, so it leaves its endomorphism
ring unchanged.  The second inequality is the Masser--W\"ustholz
estimate established in \Cref{ht:endomorphisms}; in particular its
constants have polynomial dependence on the varying field degree.
It estimates the full geometric order after an extension of bounded
relative degree, which has already been accounted for there.

For clarity, the trace conventions agree exactly in this case.
For $D_s\simeq\Mat_2(\Q)$, its action on $H_1(A_s,\Q)$ is a sum of
two standard two-dimensional modules.  Left multiplication on $D_s$
also acts on its two columns as two standard modules.  Consequently
\begin{equation}\label{fin:trace-normalization}
 \Tr_{H_1}(x)=2\tr(x)=\Tr_{D_s/\Q}(x).
\end{equation}
Thus \eqref{fin:disc-definition} is precisely the regular-trace
discriminant used in \Cref{ht:endomorphisms} and in
\cite[\S\S3.1 and 5.4]{DawOrr2022}. Bounding a smaller lattice generated by
the two symmetric endomorphisms alone would not be this assertion.

Combining \eqref{fin:degree-descent} and
\eqref{fin:height-to-disc}, and increasing constants, yields
\[
 \Delta_s\leq C_3 r_s^{v(u+1)}.
\]
Each point of the finite set $F\cap\Sigma_{E^2}(C)$ has a finite
positive discriminant.  Increasing $C_3$ includes these points too.
Taking $\delta=1/(v(u+1))$ and $c=C_3^{-\delta}$ proves the lower
bound in \eqref{fin:orbit-bound}.

Finally, for an algebraic point of a variety defined over the number
field $K$, its distinct $K$-embeddings are precisely its Galois
conjugates.  Every embedding of its residue field into $\C$ extends
to an automorphism of $\C$ over $K$.  Hence their number is
$[K(s):K]$, proving the asserted orbit equality.
\end{proof}

\subsection{All polarizations and all elliptic-square curves}

An $E^2$-special curve will mean a special curve whose generic rational
endomorphism algebra is $\Mat_2(\Q)$, as in
\cite[\S6.2]{DawOrr2022}.  We verify the geometric and integral
identifications needed to apply its theorem.

\begin{proposition}[Coverage of the square locus]\label{fin:squares}
An algebraic point $s\in\A_2(\Qbar)$ is isogenous over $\Qbar$ to
$E^2$ with $E$ non-CM if and only if it is Hodge generic on an
$E^2$-special curve.  This statement allows the given arbitrary
principal polarization on $A_s$ and an arbitrary unpolarized isogeny.
At such a point, the complexity $\Delta(Z_s)$ of its special curve
in \cite[\S6.3]{DawOrr2022} equals $\Delta_s$.
\end{proposition}

\begin{proof}
Fix an embedding into $\C$, put $V=H_1(E,\Q)$, and let $W=\Q^2$
with its trivial Hodge structure.  The given unpolarized isogeny
provides a rational Hodge isomorphism
\begin{equation}\label{fin:tensor-decomposition}
 H_1(A_s,\Q)\simeq V\otimes W.
\end{equation}
Choose an elliptic polarization $\psi_E$ on $V$.  Since $E$ is non-CM,
$\End_{\mathrm{Hdg}}(V)=\Q$.  The polarization $\psi_E$ identifies
$V$ with its appropriately twisted dual, so every Hodge pairing of
two copies of $V$ into $\Q(1)$ is a rational multiple of $\psi_E$.
Apply this to the four matrix entries, on the two copies of $V$, of
the alternating form $\psi_s$ defined by the actual polarization
$\lambda_s$.  There is a rational bilinear form $b$ on $W$ with
\begin{equation}\label{fin:polarization-form}
 \psi_s=\psi_E\otimes b.
\end{equation}
Alternation of $\psi_s$ and $\psi_E$ makes $b$ symmetric.  Let $J$
be the complex structure on $V_{\R}$, choosing the common sign
convention for which $\psi_E(v,Jv)>0$ for $v\ne0$.
The positivity of $\psi_s$ gives
\[
 0<\psi_s(v\otimes w,Jv\otimes w)
   =\psi_E(v,Jv)b(w,w)
 \quad (v\ne0,\ w\ne0),
\]
so $b$ is positive definite over $\R$.

Let $G$ be the image of
$\GL(V)=\GSp(V,\psi_E)$ acting on $V\otimes W$ through the first
factor.  In dimension two, $g$ multiplies $\psi_E$ by $\det(g)$;
therefore \eqref{fin:polarization-form} gives a rational inclusion
\[
 G\ \subset\ \GSp(V\otimes W,\psi_s).
\]
The elliptic Hodge homomorphism, acting trivially on $W$, is the
Hodge homomorphism of $A_s$.  Its $G(\R)$-conjugacy class is the
elliptic upper and lower half-plane datum.  Positivity of $b$ ensures
that this class lies in the Siegel datum for $\psi_s$.  Thus it gives
a one-dimensional rational Shimura subdatum through the period
point of $A_s$.

Retain the actual lattice $\Lambda=H_1(A_s,\Z)$ inside
$V\otimes W$.  The principal polarization makes
$(\Lambda,\psi_s)$ unimodular; a symplectic basis identifies it with
the standard integral Siegel lattice.  Intersecting its adelic
stabilizer with $G(\mathbb A_f)$ gives a compact open subgroup of
$G(\mathbb A_f)$.  The associated morphism of Shimura varieties
therefore has a special curve as the image of the component through
our point.  This construction uses $\Lambda$ itself: neither a
product lattice nor a product principal polarization has been
substituted for it.  Varying the rational identifications and the
lattices also allows every Hecke translate.

The Mumford--Tate group of a non-CM elliptic curve is $\GL(V)$.
One can see the relevant dichotomy directly: its connected Hodge
group is a reductive subgroup of $\operatorname{SL}(V)$, and a proper such subgroup
containing the elliptic Hodge circle is a torus.  That torus case
has an imaginary quadratic commutant and is precisely CM.  Therefore
the point in \eqref{fin:tensor-decomposition} has Mumford--Tate
group $G$, and is Hodge generic on the special curve just constructed.
The commutant of $G$ on $V\otimes W$ is
$\id_V\otimes\End(W)\simeq\Mat_2(\Q)$, so this is an
$E^2$-special curve.

Conversely, the defining rational representation of an
$E^2$-special curve has this same two-copy elliptic form.  At a Hodge
generic point its commutant is $\Mat_2(\Q)$ and its elliptic factor
is non-CM.  Equivalently, the rational matrix idempotents split its
abelian surface, up to isogeny, into two isogenous elliptic factors;
the off-diagonal matrix units identify their rational Hodge
structures.  The equivalence between abelian varieties up to
isogeny and polarizable rational Hodge structures of type
$\{(-1,0),(0,-1)\}$ makes this a complex isogeny
\cite[Theorem~4.1]{MilneShimura1995}.
For an algebraic point, all these homomorphisms and abelian
subvarieties descend to $\Qbar$: each is defined over a finitely
generated extension of $\Qbar$, which is regular, and
\cite[Corollary~20.4]{MilneAV} applies.  Thus the isogeny and the
elliptic factors can be taken over $\Qbar$.

It remains to check the order, rather than only its rational algebra.
For the rational subgroup $G$ defining the curve, its integral
commutant is
\begin{equation}\label{fin:integral-commutant}
 R=\End_G(H_1(A_s,\Q))\cap\End_{\Z}(\Lambda).
\end{equation}
At a Hodge generic point, commuting with $G$ is equivalent to being
a Hodge endomorphism.  To check the integral refinement, let
$u\in R$.  The rational Hodge equivalence writes $u=f/n$ for an
endomorphism $f$ of $A_s$ and a positive integer $n$.
The inclusion $u\Lambda\subset\Lambda$ means
$f\Lambda\subset n\Lambda$.  Under
$A_s[n](\C)\simeq n^{-1}\Lambda/\Lambda$, it follows that $f$
annihilates $A_s[n]$.  Factoring $f$ through the quotient by
$A_s[n]$, namely multiplication by $n$ on $A_s$, makes $u$ an
actual endomorphism.  Conversely every abelian-variety endomorphism
preserves the integral Hodge structure.
The same descent argument identifies its complex and geometric
endomorphism rings.  It follows that $R=\cO_s$.
Daw--Orr define $\Delta(Z_s)$ as the absolute discriminant of this
integral commutant \cite[\S6.3]{DawOrr2022};
\eqref{fin:trace-normalization} proves $\Delta(Z_s)=\Delta_s$.
\end{proof}

\subsection{The counting implication}

We use the following precise $E^2$ case of the published conditional
finiteness theorem \cite[Theorem~1.3, Conjecture~6.2, and
\S\S6.5--6.6]{DawOrr2022}.  Let $C\subset\A_2$ be an irreducible
Hodge-generic algebraic curve.  Let $\Sigma^{\mathrm{DO}}_{E^2}$
be the points Hodge generic on $E^2$-special curves.
If, over a finitely generated field $L\subset\C$ defining $C$,
there are $c_0,\delta_0>0$ such that
\begin{equation}\label{fin:DO-hypothesis}
 \#\bigl(\operatorname{Aut}(\C/L)\cdot s\bigr)
   \geq c_0\Delta(Z_s)^{\delta_0}
 \qquad(s\in C\cap\Sigma^{\mathrm{DO}}_{E^2}),
\end{equation}
then $C\cap\Sigma^{\mathrm{DO}}_{E^2}$ is finite.  Here $Z_s$ is
the unique special curve containing the non-special point $s$.
The constants may depend on $C$ and $L$, and no boundary condition
occurs in this implication.

\begin{proof}[Proof of \Cref{intro:squares}]
Every special curve is defined over $\Qbar$.  Since the closed
Hodge-generic curve $C$ is not contained in any such curve, its
intersection with each one is a zero-dimensional algebraic set
over $\Qbar$.  Consequently every point of
$C(\C)\cap\Sigma^{\mathrm{DO}}_{E^2}$ is algebraic.
By \Cref{fin:squares}, this intersection is exactly
$\Sigma_{E^2}(C)$, and its complexity is $\Delta_s$.
\Cref{fin:orbit} supplies \eqref{fin:DO-hypothesis} with $L=K$,
a number field and therefore finitely generated.  The stated
counting implication proves that $\Sigma_{E^2}(C)$ is finite.
In particular, in the infinite case used to choose the fixed center,
it gives the required contradiction. No companion branch theorem is used.
\end{proof}

\subsection{The special-curve classification}

The passage from the three branch statements to the full conjecture
requires a classification of the special curves, including their Hecke
translates. We record the precise consequence of the classification of
Hodge groups of abelian surfaces by Moonen and Zarhin
\cite[Theorem~(0.1)(iv), \S2 and Corollary~(3.9)]{MoonenZarhin1999}.
The resulting list of special subvarieties is also given explicitly
in \cite[\S2.F, Table~1]{DawOrrECM2021}.

\begin{proposition}[The three special-curve types]\label{fin:classification}
Let $Z\subset\A_2$ be a special curve and let $s\in Z(\Qbar)$
be Hodge generic on $Z$. Exactly one of the following holds:
\begin{enumerate}[label=\textup{(\roman*)}]
\item $A_s$ is isogenous to $E\times F$, where $E$ is non-CM
and $F$ is CM;
\item $A_s$ is isogenous to $E^2$, where $E$ is non-CM;
\item the full algebra $\End^0(A_s)$ is an indefinite quaternion
division algebra over $\Q$.
\end{enumerate}
Every point of a special curve that is not Hodge generic on it is a
special point.
\end{proposition}

\begin{proof}
The Hodge group is the smallest connected $\Q$-algebraic subgroup
whose real points contain the Hodge circle; adjoining the weight
homotheties gives the Mumford--Tate group.
At a Hodge-generic point of a special subvariety, its associated
Hermitian domain has the dimension of that subvariety. We apply this
with dimension one, retaining the actual rational Hodge structure and
polarization of $A_s$.

Suppose first that $A_s$ is simple up to isogeny. The surface
classification in \cite[\S2(2.2)]{MoonenZarhin1999} gives four
possibilities for its full rational endomorphism algebra: $\Q$, a
real quadratic field, an indefinite quaternion division algebra, or
a quartic CM field. The corresponding Hodge groups are respectively
$\Sp_4$, a restriction of scalars of $\SL_2$ from the real quadratic
field, the norm-one quaternion group, and a torus. Their Hermitian
domains have dimensions $3,2,1,0$: respectively the genus-two Siegel
space, two copies of the upper half-plane, one upper half-plane, and
a point. Thus only the quaternionic case can occur on a special curve
at a Hodge-generic point. This gives \textup{(iii)} as an equality of
the full endomorphism algebra.

If $A_s$ is not simple, Poincar\'e reducibility gives an isogeny
$A_s\sim E_1\times E_2$. When the two elliptic factors are isogenous,
their product has the same Hodge group as either factor
\cite[\S1(1.2)]{MoonenZarhin1999}. For a non-CM factor this group
is $\SL_2$, with one-dimensional domain, whereas a CM factor has
a torus Hodge group and a zero-dimensional domain. This gives
\textup{(ii)}. When the factors are nonisogenous, their Hodge group
is the product of the two elliptic Hodge groups
\cite[Corollary~(3.9)]{MoonenZarhin1999}. Each non-CM factor contributes
one domain dimension and each CM factor contributes zero. Dimension
one therefore means precisely one CM factor, giving \textup{(i)}.
The alternatives are disjoint, either by simplicity or by the full
endomorphism algebras $\Q\times F$ and $M_2(\Q)$.

These assertions are invariant under rational Hodge isomorphism, so
they apply to every polarization and Hecke translate. The elliptic
factors and isogenies at an algebraic point can be defined over
$\Qbar$, by the descent explained in \Cref{fin:squares}. Finally the
special closure of any point on $Z$ is a special subvariety contained
in $Z$. If it is proper, it has dimension zero and is a special point.
This proves the last assertion.
\end{proof}

\subsection{The full Zilber--Pink theorem}\label{fin:complete}

\begin{proof}[Proof of \Cref{intro:zp}]
First the set of special points on $C$ is finite.  These are the CM
points of $\A_2$.  If there were infinitely many, their Zariski
closure would be all of the irreducible curve $C$, since a proper
closed subset of an irreducible algebraic curve is finite.
The Andr\'e--Oort theorem for the coarse moduli space of principally
polarized abelian surfaces would then make $C$ special
\cite[Theorem~1.1]{PilaTsimerman2013}.  This contradicts its Hodge
genericity.

By \Cref{fin:classification}, after removing the already finite set
of special points, every element of $\Sigma(C)$ belongs to one of
the three indicated loci. The classification applies to all Hecke
translates and to the original coarse moduli space.

Theorem~1.1 of the $E\times\mathrm{CM}$ companion
\cite{CompanionECM}, restated as \Cref{intro:branches}(A), gives
finiteness of the entire first locus, with arbitrary CM fields,
orders, and unpolarized isogenies. \Cref{intro:squares}, proved
above, gives finiteness of the second locus. Theorem~1.1 of the
quaternionic companion \cite{CompanionQuaternion}, restated as
\Cref{intro:branches}(B), gives finiteness of the third locus:
its full endomorphism algebra is an indefinite quaternion division
algebra over $\Q$, exactly as the companion theorem requires.
The union of these three finite sets and the finite set of special
points contains $\Sigma(C)$, which is therefore finite.

All these conclusions concern the original coarse curve.  The
normalizations and fixed covers used for the second locus have
already been descended in \Cref{fin:orbit}, and every fixed omission
has been restored there.  A singular original curve has only
finitely many points removed in this process.  The argument is also
valid when $C$ is complete, because none of these steps requires
a point of its closure on a compactification boundary.
\end{proof}

The special-closure formulation of Pink's conjecture
\cite[Conjecture~1.2]{Pink2005} follows for every curve over $\Qbar$.

\begin{corollary}[The special-closure formulation]
\label{fin:special-closure}
Let $C\subset\A_2$ be a nonempty reduced irreducible closed algebraic
curve over $\Qbar$, and let $S_C$ be the smallest Shimura-special
subvariety containing $C$. Then
\[
 C(\Qbar)\cap
 \bigcup_{\substack{Z\subset\A_2\ \mathrm{special}\mathstrut\\
                     \dim Z<\dim S_C-1}} Z(\Qbar)
\]
is finite.
\end{corollary}

\begin{proof}
Write $d=\dim S_C\in\{1,2,3\}$. If $d=3$, then $C$ is Hodge generic
in $\A_2$, and the conclusion is \Cref{intro:zp}. If $d=2$, the union
consists only of special points. Infinitely many such points on $C$
would be Zariski dense, so Andr\'e--Oort
\cite[Theorem~1.1]{PilaTsimerman2013} would make $C$ special,
contradicting $\dim S_C=2$. Finally, if $d=1$, the union is empty.
\end{proof}
